\chapter{Conclusion and Future Work}
\label{sec:conclusion}

\section{Summary of findings}
\label{sec:concl-summary}

This thesis asked whether a compact language model, specialised in a single decision and used as a decision model, can
verify retrieval-augmented answers sentence by sentence on a laptop, and where any advantage comes from. GroundOne, a
fine-tuned Qwen2.5 model that reads the verdicts of a whole answer in one forward pass, gives a positive but qualified
answer.

Specialisation works where it was aimed. On the RAGTruth test set GroundOne reaches a three-way balanced accuracy of
0.769 with 0.5B parameters and 0.832 with 3B, far above the same models untrained, above the specialised encoders in
ranking supported sentences and well above Rizzo-Flow, a general-purpose decision model eight times larger. On LLM-AggreFact the 3B is on
par with the best small encoder, and on VitaminC the gain disappears: a specialist is not a general fact checker. d1-3B,
a general-purpose decision model of the same size trained to decide on a broad mixture of tasks, confirms the trade-off:
GroundOne is far ahead of it on RAGTruth, on par on LLM-AggreFact and far behind on VitaminC.

The decision-model mechanism pays off in two ways. Reading all verdicts in one pass makes verification about two times
faster than a terse and six times faster than a verbose generative judge of the same size on long answers, with output
that is always valid; and restricting the decision to the admissible answers is, by itself, more accurate than letting
the same model generate its verdict. At equal output, generative judges cannot provide probabilities of comparable
quality: their probabilities barely separate correct from incorrect decisions, and verbalised ones are almost always
all-or-nothing.

The probabilities of GroundOne can be calibrated with a temperature and a bias fitted on development data, and a
threshold chosen on the same data turns them into labels that agree with the probabilities and with the evaluated
decision. The expected calibration error alone, however, does not distinguish good probabilities from useless ones, and
discrimination must be measured separately.

The whole system runs on a 16\,GB laptop. A long answer is verified in 0.41 seconds and 5.5 joules with the 0.5B, and in
2.57 seconds and 33 joules with the 3B. The energy saved with respect to a judge is real but smaller than the time saved,
because a single long prefill draws more power than decoding.

The main limitations are the narrow training data, which leave contradictions outside data-to-text and hallucinations
hidden in copied sentences only partly learned, and a training recipe that is not fully robust to the random seed.

\section{Future work}
\label{sec:concl-future}

Several directions follow directly from these limits. The training recipe should be made robust, by enabling bias
correction and gradient clipping, logging gradient norms and repeating the runs over several seeds, including for the 3B.
Training data with more varied contradictions, and with paraphrased support and minimal edits of the kind VitaminC
contains, should address both the weak recall outside data-to-text and the blind spot on copied sentences. Training the
readout on the letter that the model would naturally write after the marker could align readout and generation, with
effects on accuracy that remain to be measured.

On the systems side, the cost model of Chapter~\ref{sec:local} should be checked on other hardware: a port of the
readout to llama.cpp would allow measuring GroundOne on a discrete graphics card, on a DGX Spark or an AMD Strix Halo
machine, and on a phone, whose neural processing units accelerate exactly the prefill on which a decision model relies. The
remaining measurement of the locality pillar is the memory of the retriever next to the generator and the verifier,
which were already measured together in this thesis. A fair latency comparison with the
encoders requires running them in batches, and one with the general-purpose decision models requires comparing both systems
through the same HTTP interface. For contexts made of several documents, a natural extension is attribution: reading, besides the verdict, which
passage supports or contradicts each sentence, for example with a readout per passage, so as to show users the source of
every claim and to handle conflicting sources explicitly. Verification could also run while the answer is being generated, flagging each sentence
as soon as it is complete, which the marker design makes natural. Finally, a local assistant built on the served
interface would let users choose the policy and the threshold that suit their documents, which is where calibrated
probabilities show their value.
